%auto-ignore

\begin{figure}[b]
\begin{tikzpicture}
 \begin{axis}[
   width=0.95\columnwidth,
   height=0.7\columnwidth,
   legend cell align=left,
   legend style={at={(1, 0)},anchor=south east,font=\scriptsize},
   mark options={mark size=3},
   font=\scriptsize,
   xmin=0, xmax=1000,
   ymin=75, ymax=85,
   xtick={200,400,600,800,1000},
   ymajorgrids=true,
   xmajorgrids=true,
   xlabel style={yshift=0.5ex,},
   xlabel={사전학습 스텝 (천 단위)} ,
   ylabel={MNLI Dev 정확도},
   ylabel style={yshift=-0.5ex,}]
    \addplot[mark=triangle,g-blue] plot coordinates {
      (30, 78.6)
      (50, 79.6)
      (100, 80.5)
      (200, 82.2)
      (400, 83.2)
      (600, 84.0)
      (800, 84.3)
      (1000, 84.4)
    };
    \addlegendentry{\bertbase (Masked LM)}
    \addplot[mark=x,g-red] plot coordinates {
      (30, 79.4)
      (50, 79.8)
      (100, 80.3)
      (200, 81.0)
      (400, 81.7)
      (600, 81.9)
      (800, 82.1)
      (1000, 82.2)
    };
    \addlegendentry{\bertbase (Left-to-Right)}
     \end{axis}
\end{tikzpicture}
\caption{학습 스텝 수에 대한 절제. $k$ 스텝 동안 사전학습된 모델 파라미터에서 출발해 파인튜닝한 뒤의 MNLI 정확도를 보여 준다. x축은 $k$의 값이다.}
\label{fig:step_abalation}
\end{figure}